% ==============================================================================
% Plantilla Institucional de Trabajo Terminal — UPIIZ IPN
% Archivo: config/comandos.tex
% ==============================================================================
% Macros, lógica de selección de tipo de documento (PROTOCOLO / TTI / TTII),
% despacho de portadas, firmas dinámicas y referencias cruzadas.
% ==============================================================================

\usepackage{etoolbox}

% ------------------------------------------------------------------------------
% 1. SISTEMA CENTRALIZADO DE TIPO DE DOCUMENTO (PROTOCOLO / TTI / TTII)
% ------------------------------------------------------------------------------
% Definición de booleanos de estado global
\newif\ifProtocolo
\newif\ifTTI
\newif\ifTTII

% Macro principal de configuración
\newcommand{\TipoDocumento}[1]{%
  \edef\tipoDocParam{\detokenize{#1}}%
  \Protocolofalse
  \TTIfalse
  \TTIIfalse
  % Comparación normalizada
  \ifstrequal{#1}{PROTOCOLO}{%
    \Protocolotrue
    \def\TipoDocCodigo{PROTOCOLO}%
  }{%
  \ifstrequal{#1}{protocolo}{%
    \Protocolotrue
    \def\TipoDocCodigo{PROTOCOLO}%
  }{%
  \ifstrequal{#1}{TTI}{%
    \TTItrue
    \def\TipoDocCodigo{TTI}%
  }{%
  \ifstrequal{#1}{tti}{%
    \TTItrue
    \def\TipoDocCodigo{TTI}%
  }{%
  \ifstrequal{#1}{TT1}{%
    \TTItrue
    \def\TipoDocCodigo{TTI}%
  }{%
  \ifstrequal{#1}{TTII}{%
    \TTIItrue
    \def\TipoDocCodigo{TTII}%
  }{%
  \ifstrequal{#1}{ttii}{%
    \TTIItrue
    \def\TipoDocCodigo{TTII}%
  }{%
  \ifstrequal{#1}{TT2}{%
    \TTIItrue
    \def\TipoDocCodigo{TTII}%
  }{%
    \PackageWarning{PlantillaUPIIZ}{%
      Tipo de documento '#1' no reconocido. Opciones validas: PROTOCOLO, TTI, TTII. %
      Se usara TTI por defecto.%
    }%
    \TTItrue
    \def\TipoDocCodigo{TTI}%
  }}}}}}}}%
}

% Macros auxiliares para consulta de estado
\newcommand{\EsProtocolo}[1]{\ifProtocolo #1\fi}
\newcommand{\EsTTI}[1]{\ifTTI #1\fi}
\newcommand{\EsTTII}[1]{\ifTTII #1\fi}

% Nombres legibles para títulos y texto
\newcommand{\NombreTipoDocumento}{%
  \ifProtocolo Protocolo de Trabajo Terminal\fi%
  \ifTTI Trabajo Terminal I\fi%
  \ifTTII Trabajo Terminal II\fi%
}

\newcommand{\NombreTipoDocumentoMayusculas}{%
  \ifProtocolo PROTOCOLO DE TRABAJO TERMINAL\fi%
  \ifTTI TRABAJO TERMINAL I\fi%
  \ifTTII TRABAJO TERMINAL II\fi%
}

\newcommand{\subtituloPortada}{%
  \ifProtocolo PROTOCOLO DE TRABAJO TERMINAL\fi%
  \ifTTI REPORTE TÉCNICO DE TRABAJO TERMINAL I\fi%
  \ifTTII REPORTE FINAL DE TRABAJO TERMINAL\fi%
}

\newcommand{\tituloValidacion}{%
  \ifTTI Análisis y validación del diseño\else Análisis y validación de resultados\fi%
}

% ------------------------------------------------------------------------------
% 2. VALIDACIÓN DE CAMPOS OBLIGATORIOS (ADVERTENCIAS EN LOG)
% ------------------------------------------------------------------------------
\newcommand{\validarDatosDocumento}{%
  \ifcsvoid{tituloProyecto}{%
    \PackageWarning{PlantillaUPIIZ}{Falta definir \protect\tituloProyecto en config/datos.tex}%
  }{}%
  \ifcsvoid{alumnoANombre}{%
    \PackageWarning{PlantillaUPIIZ}{Falta definir \protect\alumnoANombre en config/datos.tex}%
  }{}%
  \ifcsvoid{asesorANombre}{%
    \PackageWarning{PlantillaUPIIZ}{Falta definir \protect\asesorANombre en config/datos.tex}%
  }{}%
  \ifProtocolo
    \ifcsvoid{areaUbicacion}{%
      \PackageWarning{PlantillaUPIIZ}{Para PROTOCOLO se sugiere definir \protect\areaUbicacion en config/datos.tex}%
    }{}%
    \ifcsvoid{intencionTitulacion}{%
      \PackageWarning{PlantillaUPIIZ}{Para PROTOCOLO se sugiere definir \protect\intencionTitulacion en config/datos.tex}%
    }{}%
  \fi
}

% ------------------------------------------------------------------------------
% 3. COMANDOS DE IDENTIFICACIÓN DINÁMICA DE AUTORES Y ASESORES
% ------------------------------------------------------------------------------

% Imprime los autores en la Portada Principal
\newcommand{\imprimirAutoresPortada}{%
  \begin{minipage}[t]{0.9\textwidth}
    \centering\fontsize{13}{16}\selectfont
    \textbf{Presenta(n):}\\[0.5em]
    \alumnoANombre%
    \ifcsvoid{alumnoBNombre}{}{%
      \\[0.3em]\alumnoBNombre%
    }%
    \ifcsvoid{alumnoCNombre}{}{%
      \\[0.3em]\alumnoCNombre%
    }%
  \end{minipage}%
}

% Imprime los asesores en la Portada Principal
\newcommand{\imprimirAsesoresPortada}{%
  \begin{minipage}[t]{0.9\textwidth}
    \centering\fontsize{13}{16}\selectfont
    \textbf{Asesor(es):}\\[0.5em]
    \asesorANombre%
    \ifcsvoid{asesorBNombre}{}{%
      \\[0.3em]\asesorBNombre%
    }%
    \ifcsvoid{asesorCNombre}{}{%
      \\[0.3em]\asesorCNombre%
    }%
  \end{minipage}%
}

% Línea de firma para portada interna
\newcommand{\lineaFirma}[2]{%
  \parbox[t]{0.42\textwidth}{%
    \centering
    \vspace{2.5em}
    \rule{0.38\textwidth}{0.6pt}\\
    \textbf{#1}\\[0.2em]
    {\small #2}
  }%
}

% ------------------------------------------------------------------------------
% 4. COMANDOS RÁPIDOS DE REFERENCIAS CRUZADAS
% ------------------------------------------------------------------------------
\newcommand{\figref}[1]{Figura~\ref{#1}}
\newcommand{\tabref}[1]{Tabla~\ref{#1}}
\newcommand{\secref}[1]{Sección~\ref{#1}}
\newcommand{\capref}[1]{Capítulo~\ref{#1}}
\newcommand{\apref}[1]{Apéndice~\ref{#1}}
\newcommand{\eqrefcmd}[1]{Ecuación~\eqref{#1}}

% ------------------------------------------------------------------------------
% 5. COMANDOS MECATRÓNICOS Y MATEMÁTICOS COMUNES
% ------------------------------------------------------------------------------
\newcommand{\matriz}[1]{\begin{bmatrix}#1\end{bmatrix}}
\newcommand{\vectorv}[1]{\begin{pmatrix}#1\end{pmatrix}}
\newcommand{\transp}{^{\mathsf{T}}}
\newcommand{\inv}{^{-1}}
